\chapter{Sécurité, Tests et Validation}

\section*{Introduction}
\addcontentsline{toc}{section}{Introduction}

La sécurité constitue un pilier fondamental de toute plateforme SaaS, en particulier lorsqu'elle manipule des données financières sensibles (comptabilité, fiches de paie, données CNSS). Ce chapitre présente l'architecture de sécurité mise en place dans StartUpLab, les mécanismes d'authentification et d'autorisation, la stratégie de tests adoptée, ainsi que les scénarios de validation fonctionnelle exécutés pour garantir la fiabilité de la plateforme.

La plateforme StartUpLab traite plusieurs catégories de données sensibles : informations personnelles des utilisateurs, données comptables d'entreprise, fiches de paie avec numéros CNSS et CIN, contrats de travail avec signatures numériques, et informations bancaires des employés. La protection de ces données impose une approche de sécurité multicouche, conforme aux bonnes pratiques internationales et aux exigences de la législation tunisienne en matière de protection des données personnelles.

\section{Architecture de sécurité}

\subsection{Approche de sécurité multicouche}

StartUpLab implémente une architecture de sécurité en profondeur (\textit{Defense in Depth}) qui comprend plusieurs couches de protection complémentaires. La figure~\ref{fig:security-layers} illustre cette architecture.

\begin{figure}[H]
\centering
\fbox{\parbox{13cm}{
\centering
\textbf{Couche 1 : RÉSEAU \& TRANSPORT} -- HTTPS/TLS, CORS, Rate Limiting \\[4pt]
\rule{12cm}{0.4pt}\\[4pt]
\textbf{Couche 2 : AUTHENTIFICATION} -- JWT, bcrypt, Google OAuth 2.0, Face Recognition \\[4pt]
\rule{12cm}{0.4pt}\\[4pt]
\textbf{Couche 3 : AUTORISATION} -- Middleware auth, Rôles (user/admin), Plans \\[4pt]
\rule{12cm}{0.4pt}\\[4pt]
\textbf{Couche 4 : VALIDATION DES DONNÉES} -- express-validator, Requêtes préparées SQL \\[4pt]
\rule{12cm}{0.4pt}\\[4pt]
\textbf{Couche 5 : DONNÉES} -- Hashage mots de passe, Tokens temporaires, Foreign Keys
}}
\caption{Architecture de sécurité multicouche (Defense in Depth)}
\label{fig:security-layers}
\end{figure}

\subsection{Sécurité au niveau du transport}

Le serveur Express.js est configuré avec des middlewares de sécurité essentiels :

\begin{table}[H]
\centering
\caption{Mécanismes de sécurité au niveau transport}
\label{tab:transport-security}
\begin{tabular}{|l|l|l|}
\hline
\textbf{Mécanisme} & \textbf{Implémentation} & \textbf{Rôle} \\
\hline
CORS & \texttt{cors()} middleware & Contrôle des origines autorisées \\
\hline
Limite de taille & \texttt{express.json(\{limit: '10mb'\})} & Protection contre les surcharges \\
\hline
Fichiers statiques & \texttt{express.static} contrôlé & Exposition limitée du FS \\
\hline
Variables d'env. & \texttt{dotenv} & Séparation secrets / code \\
\hline
\end{tabular}
\end{table}

\subsection{Matrice de sécurité par module}

\begin{table}[H]
\centering
\caption{Matrice de sécurité par module applicatif}
\label{tab:security-matrix}
\begin{tabular}{|l|l|l|l|}
\hline
\textbf{Module} & \textbf{Données sensibles} & \textbf{Niveau} & \textbf{Mécanismes} \\
\hline
Authentification & Mots de passe, tokens & Critique & bcrypt, JWT, OAuth 2.0 \\
\hline
Comptabilité & Transactions, bilans & Élevé & Auth obligatoire, isolation userId \\
\hline
Paie \& RH & Salaires, CIN, CNSS & Critique & Auth + vérification propriétaire \\
\hline
Contrats & Signatures numériques & Élevé & Hash signature, horodatage \\
\hline
Export Expert & Données partagées & Élevé & Token temporaire, expiration \\
\hline
Admin & Toutes données & Critique & Double vérification, session 8h \\
\hline
Marketplace & Activations, paiements & Élevé & Validation plan, approbation admin \\
\hline
\end{tabular}
\end{table}

\section{Authentification et autorisation}

\subsection{Système d'authentification triple}

StartUpLab propose trois méthodes d'authentification complémentaires, chacune répondant à un cas d'usage différent.

\subsubsection{Méthode 1 : Email + Mot de passe}

Le flux d'inscription et de connexion par email suit les étapes suivantes :

\begin{enumerate}
    \item L'utilisateur fournit email, mot de passe (min. 6 caractères), prénom et nom
    \item \texttt{express-validator} vérifie le format de l'email et la longueur du mot de passe
    \item Vérification que l'email n'existe pas déjà dans la base
    \item Le mot de passe est hashé avec \texttt{bcrypt} (10 salt rounds, $\sim$100ms par hash)
    \item L'utilisateur est créé avec un abonnement \texttt{free} par défaut
    \item Un JWT est généré avec une expiration de 7 jours
\end{enumerate}

\begin{verbatim}
// Algorithme Blowfish, 10 salt rounds = 1024 iterations
const hashedPassword = await bcrypt.hash(password, 10);
// Verification (timing-safe)
const isMatch = await bcrypt.compare(password, hashedPassword);
\end{verbatim}

\subsubsection{Méthode 2 : Google OAuth 2.0 (Authorization Code Flow)}

Le flux OAuth 2.0 se déroule en six étapes :

\begin{enumerate}
    \item Frontend $\rightarrow$ Redirection Google $\rightarrow$ Code d'autorisation
    \item Frontend $\rightarrow$ \texttt{POST /api/auth/google/callback \{code, redirectUri\}}
    \item Backend $\rightarrow$ Échange code contre \texttt{access\_token} (googleapis)
    \item Backend $\rightarrow$ Récupération profil (email, nom, photo)
    \item Backend $\rightarrow$ Création/mise à jour utilisateur en BDD
    \item Backend $\rightarrow$ Génération JWT $\rightarrow$ Retour au frontend
\end{enumerate}

Pour les nouveaux utilisateurs Google, un mot de passe aléatoire de 32 octets est généré et hashé (l'utilisateur utilise toujours Google pour se connecter).

\subsubsection{Méthode 3 : Reconnaissance faciale (face-api.js)}

\textbf{Phase d'enregistrement :}
\begin{itemize}
    \item La caméra capture le visage, \texttt{face-api.js} extrait un descripteur (vecteur 128 dimensions)
    \item Le descripteur est stocké en JSON dans le champ \texttt{faceDescriptor}
\end{itemize}

\textbf{Phase de connexion :}
\begin{itemize}
    \item Un nouveau descripteur est extrait et envoyé au serveur avec l'email
    \item Le serveur calcule la \textbf{distance euclidienne} :
\end{itemize}

\begin{equation}
distance = \sqrt{\sum_{i=0}^{127}(descriptor_i - stored_i)^2}
\end{equation}

Si $distance \leq 0.6$ : Authentification réussie. Si $distance > 0.6$ : Rejet.

\begin{table}[H]
\centering
\caption{Paramètres de la reconnaissance faciale}
\label{tab:face-params}
\begin{tabular}{|l|l|l|}
\hline
\textbf{Paramètre} & \textbf{Valeur} & \textbf{Justification} \\
\hline
Dimensions descripteur & 128 & Standard SSD MobileNet V1 \\
\hline
Seuil de correspondance & 0.6 & Compromis sécurité/usabilité \\
\hline
Algorithme de distance & Euclidienne & Rapide pour les embeddings \\
\hline
\end{tabular}
\end{table}

\subsection{Gestion des tokens JWT}

\begin{table}[H]
\centering
\caption{Propriétés des tokens JWT}
\label{tab:jwt-props}
\begin{tabular}{|l|l|}
\hline
\textbf{Propriété} & \textbf{Valeur} \\
\hline
Payload utilisateur & \texttt{\{userId, email\}} \\
\hline
Payload admin & \texttt{\{userId: 0, email, isAdmin: true\}} \\
\hline
Expiration utilisateur & 7 jours \\
\hline
Expiration admin & 8 heures \\
\hline
Secret & Variable d'environnement \texttt{JWT\_SECRET} \\
\hline
Algorithme & HS256 (HMAC-SHA256) \\
\hline
\end{tabular}
\end{table}

Le fichier \texttt{middleware/auth.js} implémente deux middlewares :
\begin{itemize}
    \item \textbf{\texttt{authMiddleware}} (obligatoire) : Vérifie la présence et la validité du token Bearer. Rejette avec 401 si absent ou invalide.
    \item \textbf{\texttt{optionalAuth}} (optionnel) : Tente de décoder le token s'il est présent, mais ne bloque pas si absent.
\end{itemize}

\begin{verbatim}
const authMiddleware = (req, res, next) => {
  try {
    const authHeader = req.headers.authorization;
    if (!authHeader || !authHeader.startsWith('Bearer ')) {
      return res.status(401).json({
        message: 'Acces non autorise. Token manquant.'
      });
    }
    const token = authHeader.split(' ')[1];
    const decoded = jwt.verify(token, process.env.JWT_SECRET);
    req.user = decoded;
    next();
  } catch (error) {
    return res.status(401).json({
      message: 'Token invalide ou expire.'
    });
  }
};
\end{verbatim}

\subsection{Système d'autorisation par rôles et plans}

\textbf{Niveau 1 : Rôles (user vs admin)}

\begin{table}[H]
\centering
\caption{Routes et rôles requis}
\label{tab:routes-roles}
\begin{tabular}{|l|l|l|}
\hline
\textbf{Route} & \textbf{Rôle requis} & \textbf{Middleware} \\
\hline
\texttt{/api/auth/*} & Public/User & Aucun / authMiddleware \\
\hline
\texttt{/api/projects/*} & User & authMiddleware \\
\hline
\texttt{/api/accounting/*} & User & authMiddleware \\
\hline
\texttt{/api/hr/*} & User & authMiddleware \\
\hline
\texttt{/api/admin/*} & Admin & authMiddleware + isAdmin \\
\hline
\texttt{/api/products/admin/*} & Admin & authMiddleware + isAdmin \\
\hline
\end{tabular}
\end{table}

\textbf{Niveau 2 : Plans d'abonnement}

La hiérarchie des plans est définie comme suit :
\begin{center}
\texttt{free (0) < student (1) < startup (2) < founder (3) < enterprise (4)}
\end{center}

Le fichier \texttt{utils/subscription.js} définit 4 fonctions utilitaires :
\begin{itemize}
    \item \texttt{hasFeatureAccess(plan, feature)} : Vérifie l'accès à une fonctionnalité
    \item \texttt{getPlanLimits(plan)} : Retourne les limites (projets, idées, membres)
    \item \texttt{canCreateProject(plan, count)} : Vérifie la capacité de création
    \item \texttt{canGenerateIdea(plan, count)} : Vérifie la capacité de génération
\end{itemize}

\subsection{Isolation des données utilisateur}

Chaque requête API vérifie que l'utilisateur n'accède qu'à ses propres données :

\begin{verbatim}
-- Le parametre ? est extrait de req.user.userId (depuis le JWT)
SELECT * FROM accounting_transactions WHERE userId = ?;
SELECT * FROM employees WHERE userId = ?;
SELECT * FROM payslips WHERE userId = ?;
SELECT * FROM user_products WHERE userId = ?;
\end{verbatim}

Ce pattern est appliqué systématiquement à toutes les tables : \texttt{projects}, \texttt{ideas}, \texttt{employees}, \texttt{payslips}, \texttt{accounting\_transactions}, \texttt{user\_products}, etc.

\section{Protection des données et conformité}

\subsection{Protection contre les injections SQL}

StartUpLab utilise exclusivement des \textbf{requêtes préparées} (\textit{prepared statements}) avec Better-SQLite3 :

\begin{verbatim}
// Securise : parametres automatiquement echappes
db.prepare('SELECT * FROM users WHERE email = ?').get(email);
\end{verbatim}

\subsection{Validation des entrées}

\begin{table}[H]
\centering
\caption{Validations express-validator par route}
\label{tab:validations}
\begin{tabular}{|l|l|}
\hline
\textbf{Route} & \textbf{Validations appliquées} \\
\hline
\texttt{POST /auth/register} & Email valide, mdp $\geq$ 6 chars, prénom/nom requis \\
\hline
\texttt{POST /auth/login} & Email valide, mdp requis \\
\hline
\texttt{PUT /auth/password} & Mdp actuel requis, nouveau $\geq$ 6 chars \\
\hline
\texttt{POST /auth/face/register} & Descripteur = tableau valide \\
\hline
\texttt{POST /products/activate} & productId requis, vérification existence \\
\hline
\end{tabular}
\end{table}

\subsection{Sécurité du partage expert-comptable}

Le module d'export comptable implémente un mécanisme de partage sécurisé :
\begin{enumerate}
    \item \textbf{Token UUID v4} unique pour chaque lien de partage
    \item \textbf{Expiration configurable} par l'utilisateur
    \item \textbf{Accès lecture seule} pour l'expert-comptable
    \item \textbf{Traçabilité} : date de consultation enregistrée (\texttt{viewedAt})
    \item \textbf{Révocation} possible à tout moment
\end{enumerate}

\subsection{Signature numérique des contrats}

\begin{table}[H]
\centering
\caption{Champs de signature numérique des contrats}
\label{tab:signature}
\begin{tabular}{|l|l|l|}
\hline
\textbf{Champ} & \textbf{Type} & \textbf{Rôle} \\
\hline
\texttt{employerSignature} & TEXT (Base64) & Image signature employeur \\
\hline
\texttt{employeeSignature} & TEXT (Base64) & Image signature employé \\
\hline
\texttt{employerSignedAt} & DATETIME & Horodatage employeur \\
\hline
\texttt{employeeSignedAt} & DATETIME & Horodatage employé \\
\hline
\texttt{signatureHash} & TEXT & Hash d'intégrité du contrat \\
\hline
\end{tabular}
\end{table}

\subsection{Conformité réglementaire tunisienne}

\begin{table}[H]
\centering
\caption{Conformité réglementaire tunisienne}
\label{tab:conformite}
\begin{tabular}{|l|l|}
\hline
\textbf{Exigence} & \textbf{Implémentation} \\
\hline
Loi organique n\degre 2004-63 (protection données) & Isolation données, consentement \\
\hline
Code du travail (Art. 134-1) & Conservation fiches de paie \\
\hline
Déclaration CNSS & Taux conformes (9.18\% + 16.57\%) \\
\hline
TVA (Code de la TVA) & Taux standard 19\% \\
\hline
FEC (Art. 62 du CDPF) & Export format standard \\
\hline
\end{tabular}
\end{table}

\section{Stratégie de tests}

\subsection{Pyramide de tests}

StartUpLab adopte une stratégie de tests à plusieurs niveaux, conforme à la pyramide de tests (figure~\ref{fig:test-pyramid}).

\begin{figure}[H]
\centering
\fbox{\parbox{10cm}{
\centering
\vspace{6pt}
\small Tests E2E (Parcours utilisateur complets) \\[6pt]
\rule{8cm}{0.4pt}\\[6pt]
Tests d'intégration (API + Base de données) \\[6pt]
\rule{8cm}{0.4pt}\\[6pt]
Tests unitaires (Fonctions utilitaires, validation)
\vspace{6pt}
}}
\caption{Pyramide de tests de StartUpLab}
\label{fig:test-pyramid}
\end{figure}

\subsection{Tests unitaires}

\textbf{Tests de \texttt{subscription.js} :}

\begin{table}[H]
\centering
\caption{Tests unitaires -- subscription.js}
\label{tab:unit-tests-sub}
\begin{tabular}{|l|l|l|}
\hline
\textbf{Test} & \textbf{Entrée} & \textbf{Résultat attendu} \\
\hline
\texttt{hasFeatureAccess('free', 'ideaGenerator')} & Plan free & \texttt{true} \\
\hline
\texttt{hasFeatureAccess('free', 'businessPlanPdf')} & Plan free & \texttt{false} \\
\hline
\texttt{hasFeatureAccess('enterprise', 'apiAccess')} & Plan enterprise & \texttt{true} \\
\hline
\texttt{canCreateProject('free', 0)} & 0 projets & \texttt{true} \\
\hline
\texttt{canCreateProject('free', 1)} & 1 projet & \texttt{false} (limite=1) \\
\hline
\texttt{canCreateProject('pro', 100)} & 100 projets & \texttt{true} (illimité) \\
\hline
\end{tabular}
\end{table}

\textbf{Tests bcrypt :}

\begin{table}[H]
\centering
\caption{Tests unitaires -- bcrypt}
\label{tab:unit-tests-bcrypt}
\begin{tabular}{|l|l|}
\hline
\textbf{Test} & \textbf{Résultat attendu} \\
\hline
\texttt{bcrypt.hash('test123', 10)} & Hash 60 chars, préfixe \texttt{\$2a\$} \\
\hline
\texttt{bcrypt.compare('test123', hash)} & \texttt{true} \\
\hline
\texttt{bcrypt.compare('wrong', hash)} & \texttt{false} \\
\hline
Deux hash du même mot de passe & Hash $\neq$ (sel aléatoire) \\
\hline
\end{tabular}
\end{table}

\textbf{Tests JWT :}

\begin{table}[H]
\centering
\caption{Tests unitaires -- JWT}
\label{tab:unit-tests-jwt}
\begin{tabular}{|l|l|}
\hline
\textbf{Test} & \textbf{Résultat attendu} \\
\hline
\texttt{jwt.sign(\{userId:1\}, secret, \{expiresIn:'7d'\})} & Token 3 parties \\
\hline
\texttt{jwt.verify(token, secret)} & Payload décodé \\
\hline
Token expiré (\texttt{expiresIn:'0s'}) & \texttt{TokenExpiredError} \\
\hline
Mauvais secret & \texttt{JsonWebTokenError} \\
\hline
\end{tabular}
\end{table}

\subsection{Tests d'intégration API}

\textbf{Module Authentification :}

\begin{table}[H]
\centering
\caption{Tests d'intégration -- Authentification}
\label{tab:integration-auth}
\begin{tabular}{|l|l|l|l|}
\hline
\textbf{\#} & \textbf{Endpoint} & \textbf{Scénario} & \textbf{Code} \\
\hline
T1 & \texttt{POST /auth/register} & Inscription valide & 201 \\
\hline
T2 & \texttt{POST /auth/register} & Email déjà utilisé & 400 \\
\hline
T3 & \texttt{POST /auth/register} & Email invalide & 400 \\
\hline
T4 & \texttt{POST /auth/login} & Connexion valide & 200 \\
\hline
T5 & \texttt{POST /auth/login} & Mauvais mdp & 401 \\
\hline
T6 & \texttt{GET /auth/me} & Token valide & 200 \\
\hline
T7 & \texttt{GET /auth/me} & Sans token & 401 \\
\hline
T8 & \texttt{GET /auth/me} & Token expiré & 401 \\
\hline
\end{tabular}
\end{table}

\textbf{Module Marketplace :}

\begin{table}[H]
\centering
\caption{Tests d'intégration -- Marketplace}
\label{tab:integration-marketplace}
\begin{tabular}{|l|l|l|l|}
\hline
\textbf{\#} & \textbf{Endpoint} & \textbf{Scénario} & \textbf{Code} \\
\hline
T9 & \texttt{GET /products/categories} & Liste catégories & 200 \\
\hline
T10 & \texttt{GET /products/product/:slug} & Produit existant & 200 \\
\hline
T11 & \texttt{GET /products/product/inexistant} & Slug invalide & 404 \\
\hline
T12 & \texttt{POST /products/activate} & Plan suffisant & 201 \\
\hline
T13 & \texttt{POST /products/activate} & Plan insuffisant & 403 \\
\hline
T14 & \texttt{POST /products/activate} & Déjà activé & 400 \\
\hline
T15 & \texttt{POST /products/admin/approve/:id} & Admin approuve & 200 \\
\hline
T16 & \texttt{POST /products/admin/approve/:id} & Non-admin & 403 \\
\hline
\end{tabular}
\end{table}

\textbf{Module Comptabilité et RH :}

\begin{table}[H]
\centering
\caption{Tests d'intégration -- Comptabilité et RH}
\label{tab:integration-compta}
\begin{tabular}{|l|l|l|l|}
\hline
\textbf{\#} & \textbf{Endpoint} & \textbf{Scénario} & \textbf{Code} \\
\hline
T17 & \texttt{POST /accounting/transactions} & Ajout transaction & 201 \\
\hline
T18 & \texttt{GET /accounting/bilan} & Calcul bilan & 200 \\
\hline
T19 & \texttt{GET /accounting/tva} & Déclaration TVA & 200 \\
\hline
T20 & \texttt{POST /hr/employees} & Ajout employé & 201 \\
\hline
T21 & \texttt{POST /hr/payslips/generate} & Fiche de paie & 201 \\
\hline
T22 & \texttt{GET /hr/cnss/declaration} & Déclaration CNSS & 200 \\
\hline
T23 & \texttt{POST /hr/leaves} & Demande congé & 201 \\
\hline
T24 & \texttt{POST /hr/contracts} & Création contrat & 201 \\
\hline
\end{tabular}
\end{table}

\section{Tests fonctionnels et scénarios de validation}

\subsection{Scénarios End-to-End}

\textbf{Scénario 1 : Parcours complet d'un nouvel entrepreneur}

\begin{table}[H]
\centering
\caption{Scénario E2E -- Parcours entrepreneur}
\label{tab:e2e-entrepreneur}
\begin{tabular}{|l|l|l|c|}
\hline
\textbf{Étape} & \textbf{Action} & \textbf{Résultat attendu} & \checkmark \\
\hline
1 & Accéder à la landing page & Page d'accueil affichée & \checkmark \\
\hline
2 & Cliquer Commencer gratuitement & Redirection vers /register & \checkmark \\
\hline
3 & Remplir le formulaire & Validation temps réel & \checkmark \\
\hline
4 & Soumettre l'inscription & Redirection vers /dashboard & \checkmark \\
\hline
5 & Créer un nouveau projet & Projet visible dans la liste & \checkmark \\
\hline
6 & Générer une idée & Idée affichée avec score & \checkmark \\
\hline
7 & Créer un BMC & BMC sauvegardé & \checkmark \\
\hline
8 & Accéder au Pitch Deck & Éditeur fonctionnel & \checkmark \\
\hline
9 & Inviter un membre & Email envoyé & \checkmark \\
\hline
10 & Se déconnecter & Retour landing page & \checkmark \\
\hline
\end{tabular}
\end{table}

\textbf{Scénario 2 : Gestion comptable complète}

\begin{table}[H]
\centering
\caption{Scénario E2E -- Gestion comptable}
\label{tab:e2e-compta}
\begin{tabular}{|l|l|l|c|}
\hline
\textbf{Étape} & \textbf{Action} & \textbf{Résultat attendu} & \checkmark \\
\hline
1 & Activer module Comptabilité & Demande envoyée (pending) & \checkmark \\
\hline
2 & Admin approuve & Module accessible sidebar & \checkmark \\
\hline
3 & Ajouter revenu (5000 TND) & Solde mis à jour & \checkmark \\
\hline
4 & Ajouter dépense (1500 TND) & Graphique mis à jour & \checkmark \\
\hline
5 & Consulter le bilan & Actif = 3500 TND & \checkmark \\
\hline
6 & Créer déclaration TVA & TVA = 950 TND (19\%) & \checkmark \\
\hline
7 & Exporter pour expert & Lien de partage généré & \checkmark \\
\hline
8 & Expert accède au lien & Lecture seule & \checkmark \\
\hline
\end{tabular}
\end{table}

\textbf{Scénario 3 : Gestion RH et paie}

\begin{table}[H]
\centering
\caption{Scénario E2E -- Gestion RH et paie}
\label{tab:e2e-rh}
\begin{tabular}{|l|l|l|c|}
\hline
\textbf{Étape} & \textbf{Action} & \textbf{Résultat attendu} & \checkmark \\
\hline
1 & Créer employé (CDI, 2000 TND) & Fiche créée & \checkmark \\
\hline
2 & Générer fiche de paie mai & Calculs automatiques & \checkmark \\
\hline
3 & Vérifier CNSS employé & $2000 \times 9.18\% = 183.60$ TND & \checkmark \\
\hline
4 & Vérifier CNSS employeur & $2000 \times 16.57\% = 331.40$ TND & \checkmark \\
\hline
5 & Vérifier le salaire net & Brut - CNSS - IRPP = correct & \checkmark \\
\hline
6 & Demande congé (5j) & Solde : $24 \rightarrow 19$ jours & \checkmark \\
\hline
7 & Créer contrat CDI & Mode brouillon & \checkmark \\
\hline
8 & Signer (empl. + employé) & Statut : signé complet & \checkmark \\
\hline
\end{tabular}
\end{table}

\textbf{Scénario 4 : Workflow activation marketplace}

\begin{table}[H]
\centering
\caption{Scénario E2E -- Workflow marketplace}
\label{tab:e2e-marketplace}
\begin{tabular}{|l|l|l|c|}
\hline
\textbf{Étape} & \textbf{Action} & \textbf{Résultat attendu} & \checkmark \\
\hline
1 & Consulter Produits \& Solutions & 6 catégories, 30+ produits & \checkmark \\
\hline
2 & Cliquer Demo Analyse Marché & Page démo interactive & \checkmark \\
\hline
3 & Cliquer Activer & Toast En attente & \checkmark \\
\hline
4 & Admin : Offres \& Produits & Demande visible pending & \checkmark \\
\hline
5 & Admin approuve & Notification utilisateur & \checkmark \\
\hline
6 & Sidebar utilisateur & Section IA Business apparaît & \checkmark \\
\hline
7 & Cliquer dans sidebar & Page /produit/actif/ & \checkmark \\
\hline
8 & Accéder à la fonctionnalité & Page active (pas démo) & \checkmark \\
\hline
\end{tabular}
\end{table}

\subsection{Tests de régression}

\begin{table}[H]
\centering
\caption{Tests de régression}
\label{tab:regression}
\begin{tabular}{|l|l|l|}
\hline
\textbf{ID} & \textbf{Test de régression} & \textbf{Criticité} \\
\hline
R1 & Inscription + Connexion + Déconnexion & Critique \\
\hline
R2 & CRUD Projets & Haute \\
\hline
R3 & Calcul automatique TVA 19\% & Haute \\
\hline
R4 & Calcul fiche de paie (CNSS + IRPP) & Critique \\
\hline
R5 & Activation produit $\rightarrow$ apparition sidebar & Haute \\
\hline
R6 & Navigation toutes les pages (pas de page blanche) & Haute \\
\hline
R7 & Redirection Accéder à la fonctionnalité & Moyenne \\
\hline
R8 & Admin : approuver/refuser activations & Haute \\
\hline
R9 & Export expert-comptable avec token & Moyenne \\
\hline
R10 & Notifications après actions & Moyenne \\
\hline
\end{tabular}
\end{table}

\section{Gestion des erreurs et résilience}

\subsection{Stratégie à trois niveaux}

\textbf{Niveau 1 -- Prévention (validation entrées) :}

\begin{verbatim}
const errors = validationResult(req);
if (!errors.isEmpty()) {
  return res.status(400).json({ errors: errors.array() });
}
\end{verbatim}

\textbf{Niveau 2 -- Détection (logique métier) :}

\begin{verbatim}
if (userLevel < requiredLevel) {
  return res.status(403).json({
    message: `Ce produit necessite un abonnement ${product.requiredPlan}`,
    requiredPlan: product.requiredPlan
  });
}
\end{verbatim}

\textbf{Niveau 3 -- Récupération (try-catch global) :}

\begin{verbatim}
try {
  // logique metier
} catch (error) {
  console.error('Operation error:', error);
  res.status(500).json({ message: 'Erreur serveur' });
}
\end{verbatim}

\subsection{Codes HTTP utilisés}

\begin{table}[H]
\centering
\caption{Codes HTTP utilisés dans StartUpLab}
\label{tab:http-codes}
\begin{tabular}{|l|l|l|}
\hline
\textbf{Code} & \textbf{Signification} & \textbf{Usage} \\
\hline
200 & OK & Succès lecture/mise à jour \\
\hline
201 & Created & Succès création \\
\hline
400 & Bad Request & Données invalides, doublon \\
\hline
401 & Unauthorized & Token manquant/invalide \\
\hline
403 & Forbidden & Plan insuffisant, accès admin requis \\
\hline
404 & Not Found & Ressource inexistante \\
\hline
500 & Internal Server Error & Erreur technique \\
\hline
\end{tabular}
\end{table}

\subsection{Gestion frontend (react-hot-toast)}

\begin{table}[H]
\centering
\caption{Types de notifications frontend}
\label{tab:toast}
\begin{tabular}{|l|l|l|}
\hline
\textbf{Type} & \textbf{Couleur} & \textbf{Durée} \\
\hline
Succès & Vert \texttt{\#10b981} & 4 secondes \\
\hline
Erreur & Rouge \texttt{\#ef4444} & 4 secondes \\
\hline
Info & Sombre \texttt{\#363636} & 4 secondes \\
\hline
\end{tabular}
\end{table}

Un intercepteur Axios déclenche la déconnexion automatique si le serveur retourne 401 (token expiré).

\subsection{Résilience de la base de données}

\begin{table}[H]
\centering
\caption{Propriétés de résilience SQLite}
\label{tab:resilience}
\begin{tabular}{|l|l|}
\hline
\textbf{Propriété} & \textbf{Détail} \\
\hline
ACID & Transactions atomiques, cohérentes, isolées, durables \\
\hline
Foreign Keys & \texttt{PRAGMA foreign\_keys = ON} \\
\hline
Contraintes UNIQUE & Email, slugs, combinaisons userId+productId \\
\hline
CASCADE DELETE & Suppression automatique données liées \\
\hline
DEFAULT VALUES & Valeurs par défaut sûres \\
\hline
\end{tabular}
\end{table}

\section{Audit de sécurité et recommandations}

\subsection{Résultats de l'audit}

\begin{table}[H]
\centering
\caption{Résultats de l'audit de sécurité}
\label{tab:audit}
\begin{tabular}{|l|l|l|}
\hline
\textbf{Critère} & \textbf{Statut} & \textbf{Score} \\
\hline
Hashage mots de passe & Conforme & 10/10 \\
\hline
Injection SQL & Protégé & 10/10 \\
\hline
Authentification JWT & Implémenté & 9/10 \\
\hline
Contrôle d'accès rôle & Implémenté & 9/10 \\
\hline
Validation entrées & express-validator & 8/10 \\
\hline
CORS & Configuré & 8/10 \\
\hline
Variables d'environnement & dotenv & 9/10 \\
\hline
HTTPS & Dev uniquement (HTTP) & 5/10 \\
\hline
Rate limiting & Non implémenté & 3/10 \\
\hline
CSRF & SPA + JWT (immunisé) & 7/10 \\
\hline
Logs de sécurité & Console uniquement & 5/10 \\
\hline
\textbf{Score global} & & \textbf{81/110 (74\%)} \\
\hline
\end{tabular}
\end{table}

\subsection{Recommandations}

\begin{table}[H]
\centering
\caption{Recommandations de sécurité}
\label{tab:recommandations}
\begin{tabular}{|l|l|l|l|}
\hline
\textbf{Priorité} & \textbf{Recommandation} & \textbf{Effort} & \textbf{Impact} \\
\hline
Critique & HTTPS/TLS obligatoire & Faible & Élevé \\
\hline
Haute & Rate limiting (express-rate-limit) & Faible & Élevé \\
\hline
Haute & Logs structurés (Winston) & Moyen & Moyen \\
\hline
Moyenne & Headers sécurité (HSTS, CSP) & Faible & Moyen \\
\hline
Moyenne & Refresh token (JWT rotation) & Moyen & Moyen \\
\hline
Basse & 2FA & Élevé & Élevé \\
\hline
Basse & Chiffrement données au repos & Élevé & Élevé \\
\hline
\end{tabular}
\end{table}

\section*{Conclusion}
\addcontentsline{toc}{section}{Conclusion}

L'architecture de sécurité de StartUpLab couvre les menaces principales d'une plateforme SaaS. L'authentification triple, les tokens JWT, l'isolation des données par utilisateur, et la validation systématique des entrées constituent un socle solide. Les tests fonctionnels couvrent les scénarios critiques de bout en bout. L'audit révèle un score de 74\% avec des recommandations claires pour atteindre un niveau de sécurité production.
